\chapter{Cellular Respiration: Electrons Drive ATP Synthesis}

\begin{kaichang}
Stand up now and run briskly in place for a minute, or do thirty squats. Stop and notice your body: you are breathing hard, your face feels hot, you may be sweating, and your heart is pounding.

Good. Sit down, and let us ask a question while you are in that state. What ultimately happened to all the oxygen you just inhaled?

Guess first. Most people's immediate answer is, ``It became carbon dioxide, which I breathed out.'' After all, that is what physical-education lessons say: breathe oxygen in, carbon dioxide out. Write down your guess. We will check it at the end. The right answer surprises many people: most of the oxygen you inhale becomes \textbf{water}. The oxygen in carbon dioxide comes from elsewhere.

Behind this whole sequence, from panting to sweating, is a power station hidden in almost every cell. This chapter opens it up to see how electrons release energy step by step, like descending stairs, and how they drive ATP synthesis.
\end{kaichang}

\section{Combustion Without Flames}

Start with some observations anyone can make.

First, you cannot stop breathing. You can go days without food or a day without water, but complete oxygen deprivation begins irreversibly damaging the brain within minutes. Some bodily process clearly needs oxygen \textbf{continuously}; the supply cannot simply be stored for later.

Second, harder exercise means faster breathing. Sitting quietly, you breathe a dozen or more times per minute; after sprinting, that can rise to thirty or forty. The body is like an engine that shifts gears automatically: more work, more intake.

Third, exercise warms you. Run a couple of laps in winter, and you feel warmer. This is real: your temperature-control system works hard to shed heat by dilating skin blood vessels, reddening your face, and producing sweat.

Fourth, exhaled air is moist. Breathe on a cold window, and it immediately mists over. You might say that the water was simply drunk and then expelled. But consider this: even if an adult drinks nothing for a whole day, water vapor continues leaving with every breath.

Fifth, a precise measurement: burning 1 gram of glucose in a bomb calorimeter releases the same total heat as an animal completely using 1 gram of glucose, about \ud{16}{kJ}. Different paths, identical endpoints: the energy difference is exactly the same. Yet burning glucose produces flames, whereas your body has none. A flame releases the energy all at once; your body divides it into dozens of small payments and stores about a third in the ATP savings card introduced in Chapter 50.

How does the payment system work? We must zoom inside a cell.

\begin{shiyan}
\textbf{Materials:} a small mirror or cold glass surface, such as a chilled drink bottle, and a stopwatch or phone timer.

\textbf{Procedure:} (1) Sit quietly for five minutes, then count and record your breaths over one minute. (2) Run briskly in place for a minute, then immediately count and record breaths for another minute. (3) Breathe out slowly close to the mirror and observe its surface.

\textbf{Observations:} by what factor did your breathing rate change? Where did the mist on the mirror come from? Not all exhaled water was simply drunk and passing through: by the end of this chapter, you will know that some was freshly manufactured in your cells.

\textbf{Safety:} exercise within your abilities and stop if dizzy. Other material in this chapter involves poisons such as mitochondrial inhibitors. \textbf{Never} attempt related procedures at home; watch only a teacher's demonstration or a reputable video.

\textbf{Translator safety note:} References to poison demonstrations below are explanatory, not activities for students or casual classroom use. Use diagrams or reputable recordings, not cyanide or uncoupling chemicals. Suspected cyanide exposure requires immediate emergency help; see \href{https://www.cdc.gov/chemical-emergencies/chemical-fact-sheets/cyanide.html}{CDC cyanide guidance}.
\end{shiyan}

\section{Pyruvic Acid Enters the Power Station}

As Chapter 52 explained, glycolysis in the cytoplasm splits glucose into two pyruvic-acid molecules, earns a net 2 ATP, and stores some high-energy electrons in NADH. The bargain is poor: most of glucose's energy remains in pyruvic acid, as though you opened a package and took only the free gift, leaving the main item untouched.

Where is that main item dismantled? In the \textbf{mitochondrion}. First learn its structure, because this chapter's machinery sits on its membranes. A mitochondrion has \textbf{two membranes}. The outer membrane is relatively smooth, like a factory wall. The inner one repeatedly folds inward into ridges called cristae, like radiator fins, packing as much membrane area as possible into a limited volume. The inner compartment is the \textbf{matrix}; the narrow space between membranes is the \textbf{intermembrane space}. Remember these three place names---matrix, inner membrane, intermembrane space---because they recur throughout the chapter.

Incidentally, mitochondria are not distributed one per cell. Heavy workloads call for more power stations. Heart-muscle cells beat throughout your life, and mitochondria occupy a third of their volume. A ring of tightly packed mitochondria in the sperm tail powers its swimming strokes. Mature red blood cells have none: they reserve their space for hemoglobin and get by on glycolysis's 2 ATP. The distribution of your body's power stations is a precise map of energy use.

After crossing both membranes into the matrix, pyruvic acid undergoes an entrance check: one carbon is removed as \chem{CO_2}. Notice that this is one birthplace of the carbon dioxide you exhale. The remaining two-carbon fragment attaches to a carrier, coenzyme A, and enters a circular assembly line: the \textbf{citric-acid cycle}, also called the tricarboxylic-acid or Krebs cycle.

The design is clever: instead of a line with a start and finish, it is a \textbf{circle}. The two-carbon fragment joins a four-carbon molecule to form six-carbon citric acid. During a turn, two carbons leave as \chem{CO_2}, and the original four-carbon molecule is regenerated at the end to welcome the next shipment. Thus citric acid is both the first product and cyclically regenerated, giving the cycle its name.

The main harvest per turn is not ATP---only 1 ATP equivalent is obtained directly---but \textbf{electrons}. Four operations strip high-energy electrons from the carbon skeleton and load them into carriers: 3 NADH and 1 \chem{FADH_2}, a carrier like NADH among Chapter 50's electron cash vans. Add the 1 NADH produced at the entrance check, and complete breakdown of one pyruvic-acid molecule yields 4 NADH, 1 \chem{FADH_2}, and 3 \chem{CO_2}.

One glucose supplies two pyruvic-acid molecules, so double those numbers. Glucose's carbon skeleton has now \textbf{entirely} left as \chem{CO_2}, but the cell has gained only 4 ATP. Where is the energy? Most is now in the high-energy electrons carried by NADH and \chem{FADH_2}. The carbon has left; the electrons remain. Actual electricity generation is only beginning.

\section{Electrons Downstairs, Protons Uphill}

High-energy electrons are cashed in by a series of giant protein machines embedded in the mitochondrial \textbf{inner membrane}: the \textbf{electron-transport chain}. It consists mainly of four protein complexes numbered I to IV, like descending steps.

NADH hands its electrons to the first step, complex I, and they then pass to III and IV. Why stairs rather than a single drop? This is Chapter 27's energy rule applied to biology. Each descent from a higher to a lower energy level releases a small portion of free energy, giving the protein machine a chance to \textbf{capture and use it}. A single leap would dissipate everything as heat, like dropping a bucket from a tenth-floor window: all its kinetic energy becomes a loud crash. Stepwise transfer is the cell's version of water successively turning a series of small waterwheels.

What do the steps do with the captured energy? Surprisingly, they \textbf{move protons}. Complexes I, III, and IV use energy released during electron transfer to pump protons, \chem{H^+}, across the inner membrane from the matrix to the intermembrane space. Electrons go downstairs while protons go uphill. The electrons' released energy becomes the protons' height difference.

This creates a large imbalance: the intermembrane space has more protons and positive charge, while the matrix has fewer protons. Chapter 50 explained that both concentration differences and charge differences store energy. Together, they form the \textbf{proton-motive force}, like the water-level difference across a reservoir dam. It is only about \ud{0.2}{V}, but across a membrane just a few nanometers thick, the electric field reaches hundreds of thousands of volts per centimeter, stronger than the field around high-voltage transmission lines.

A dam needs a turbine to generate power. The mitochondrial turbine is a fifth protein machine, \textbf{ATP synthase}, also embedded in the inner membrane. It is the smallest rotary motor currently known. Protons stream through its channel from the intermembrane space back into the matrix, driving a rotor at over a hundred revolutions per second. The rotor's twist changes catalytic sites at the other end, pressing ADP and phosphate into ATP. Chapter 50 reminded us that ATP's energy is not in a particular high-energy bond, but in its more stable hydrolysis products and concentrations far from equilibrium. Here is the other side: ATP is \textbf{made by the force of a proton flood}. As long as the gradient persists, the motor keeps turning, and one ATP synthase can make hundreds of ATP molecules per second.

One final piece remains: where do the electrons go after descending? Someone must accept them or the chain soon jams. Waiting at the bottom is the \textbf{oxygen} you breathe. At complex IV, oxygen accepts electrons, takes two protons from the matrix, and forms water:
\[\chem{O_2} + 4\chem{H^+} + 4e^- \rightarrow 2\chem{H_2O}\]
Oxygen is the \textbf{terminal electron acceptor}, an always-open electron exit. It is neither fuel nor a direct partner of glucose; it simply accepts spent electrons at the end. That is why minutes without oxygen can be fatal. Close the exit, and the whole electron line jams; proton pumps stop, the gradient dissipates, and ATP synthase stops. The body's power stations shut down together.

\section{Writing the Whole Process in Symbols}

Now symbols enter. Start with the overall account: completely oxidizing one glucose gives the same net equation as familiar combustion:
\[\chem{C_6H_{12}O_6} + 6\chem{O_2} \rightarrow 6\chem{CO_2} + 6\chem{H_2O}\]
But this records only starting and ending points, omitting dozens of intermediate steps. Combustion drops straight to the bottom; cellular respiration divides the same energy into dozens of payments. Identical equations, radically different paths: thermodynamics sets direction, while enzymes set routes and rates. A rule emphasized throughout Volume I plays out in life.

The electron cash van's handoff is:
\[\chem{NADH} \rightarrow \chem{NAD^+} + \chem{H^+} + 2e^-\]
Unloaded \chem{NAD^+} returns to glycolysis and the citric-acid cycle to collect another shipment. This is why Chapter 52 said cells must continually regenerate \chem{NAD^+}, or even glycolysis stops.

The electron-transport chain's energy landscape can be drawn as a staircase:

\begin{center}
\begin{tikzpicture}
\draw (0,3.4) rectangle (2,4.2); \node at (1,3.8) {\small Complex I};
\draw (2.6,2.5) rectangle (4.6,3.3); \node at (3.6,2.9) {\small Complex III};
\draw (5.2,1.6) rectangle (7.2,2.4); \node at (6.2,2.0) {\small Complex IV};
\draw (7.8,0.6) rectangle (9.8,1.4); \node at (8.8,1.0) {\small \chem{O_2}$\rightarrow$\chem{H_2O}};
\draw[->,thick] (-0.2,4.6) -- (0.6,4.2); \node at (0.4,4.8) {\small NADH enters};
\draw[->,thick] (2,3.3) -- (2.6,3.0);
\draw[->,thick] (4.6,2.4) -- (5.2,2.1);
\draw[->,thick] (7.2,1.5) -- (7.8,1.1);
\draw[->,thick] (1,3.4) -- (1,2.6) node[below]{\small Pumps \chem{H^+}};
\draw[->,thick] (3.6,2.5) -- (3.6,1.7) node[below]{\small Pumps \chem{H^+}};
\draw[->,thick] (6.2,1.6) -- (6.2,0.8) node[below]{\small Pumps \chem{H^+}};
\end{tikzpicture}
\end{center}

The diagram needs a reminder: this is a human representation. Real complexes are folded masses of protein; their sizes, spacing, and colors here do not show their true appearance. The picture conveys only that electrons descend in steps, each driving a proton pump.

\begin{qiaoliang}
Add up the accounts. One glucose's full electron harvest is 2 NADH from glycolysis, 2 from pyruvic acid's entry, and 6 NADH plus 2 \chem{FADH_2} from the citric-acid cycle. The chain's exchange rates are about 2.5 ATP per NADH and 1.5 per \chem{FADH_2}. The latter enters at the second step, missing one round of proton pumping, so its rate is lower. Thus $10\times 2.5 + 2\times 1.5 = 28$ ATP come from electron transport. Add the 4 obtained directly from glycolysis and the cycle for a total of about 30--32 ATP. Without oxygen, only glycolysis's 2 are available. Aerobic respiration increases energy extraction roughly fifteenfold, a chemical prerequisite for large multicellular organisms like you to afford a brain.

Another order-of-magnitude estimate shows ATP's frantic turnover. Even lying still all day, an adult synthesizes and consumes about 40--70 kilograms of ATP daily, close to body weight, yet has only tens of grams present at any instant. ATP is not warehouse inventory but cash printed and spent every second on an assembly line. Even sitting quietly, your army of ATP synthases produces hundreds of trillions of ATP molecules every second.
\end{qiaoliang}

\section{An Idea Ridiculed for Over a Decade}

\begin{zhengju}
Today, saying ``a proton gradient drives ATP synthesis'' sounds like ordinary common sense. In 1961, mainstream researchers treated the idea as bizarre.

The prevailing explanation was the \textbf{chemical-coupling hypothesis}. Since glycolysis transfers energy through high-energy chemicals step by step, perhaps the electron-transport chain first makes some mysterious high-energy intermediate, dubbed something like X\ensuremath{\sim}P, which then makes ATP. This was reasonable analogical reasoning with a clear task: isolate the intermediate. Laboratories worldwide searched for more than a decade and found nothing. The persistent absence of a theory's central predicted substance should itself raise concern.

In 1961, the British biochemist Peter Mitchell proposed a very different \textbf{chemiosmotic hypothesis}. There was no high-energy intermediate. The electron-transport chain's real product was protons pumped across a membrane: \textbf{the proton gradient itself stored energy}. As protons flowed back through ATP synthase, that gradient could become ATP. The hypothesis gave membranes a new role, not merely bags holding things, but dams and reservoirs.

The response bordered on hostility. Textbooks then said energy could be stored only in chemical bonds, and Mitchell's theory was mocked as membrane mysticism; colleagues openly rejected it at meetings. His circumstances made things harder: poor health and tight funding eventually led him to convert a country manor into a private laboratory, the Glynn Research Institute, supported by a cattle farm and family savings.

Mitchell did not win through eloquence. He designed experiments \textbf{that could condemn his theory}. One key argument was this: if chemical coupling were right, damaging membrane integrity should at most reduce efficiency. If chemiosmosis were right, a proton-leaky membrane should completely stop ATP synthesis while electron transfer continued. The results were clear: a leaky membrane left electron transfer working but stopped ATP production. Chemical coupling could not naturally explain this.

Still more elegant evidence arrived in 1974. Stoeckenius and Racker assembled a bacterial light-driven proton pump, bacteriorhodopsin, and ATP synthase in the \textbf{same} artificial lipid vesicle, with no cell extract inside. Light activated the pump, protons moved, and ATP appeared. With nothing but pump, membrane, gradient, and ATP synthase, the supposed high-energy chemical intermediate had nowhere to hide. An old puzzle also fell into place: substances such as 2,4-dinitrophenol allow respiration to continue but produce heat instead of ATP. Chemiosmosis explained them as proton smugglers that carry protons across the membrane, dissipating the gradient so all the energy becomes heat.

Mitchell alone received the 1978 Nobel Prize in Chemistry. The lesson is not that persistence always wins; scientific history has plenty of people persisting in error. It is about \textbf{what earns a reversal}. Even a minority theory should win if it predicts leaky-membrane and reconstituted-vesicle results more accurately than its rival. Being mocked for more than a decade before evidence arrives is a familiar feature of scientific communities, not a conspiracy.
\end{zhengju}

\section{When This Model Is Not Enough}

Chemiosmosis is powerful, but understanding it requires stating its limits.

First, mitochondria are not essential. Bacteria lack them and instead put electron-transport chains and ATP synthase in their \textbf{cell membranes}, still generating power from proton gradients. Chemiosmosis is more fundamental than the organelle. Indeed, mitochondria originated from ancient bacteria, a story for later chapters.

Second, proton gradients do more than make ATP. Many bacteria use proton flow directly to rotate flagella, like one reservoir supplying both electricity and water. Some membrane transport also taps the gradient. A gradient is mainly, not exclusively, an ATP battery.

Third, electron transport and ATP synthesis are not welded together. Brown-fat mitochondria have an uncoupling protein in their inner membrane that deliberately gives protons a bypass, turning the gradient directly into heat. Newborns and hibernating animals use it to stay warm. The claim that respiration exists only to make ATP has a counterexample between your shoulder blades.

Fourth, oxygen is not the only electron exit. Many microorganisms use nitrate, sulfate, or even iron ions as terminal electron acceptors, called anaerobic respiration. They live perfectly well, but with a smaller electron-energy drop and lower yield. Oxygen is a champion acceptor because its desire for electrons---electronegativity in Chapter 11---is among the strongest, providing the longest downhill route.

Fifth, 30--32 ATP is approximate. Membranes always leak a little; transporting NADH into and out of mitochondria incurs fees; actual yield varies with cell type and state. Quantitative biology nearly always carries this word ``about.''

\begin{wuqu}
``The oxygen you inhale becomes the carbon dioxide you exhale.'' This widespread mistake is instructive. \chem{CO_2} is removed from \textbf{food's carbon skeletons} at pyruvic-acid entry and in the citric-acid cycle. Its oxygen atoms come from glucose and water, not directly from inhaled \chem{O_2}. Inhaled \chem{O_2} accepts electrons and protons at the end of the electron-transport chain and becomes \textbf{water}. The clearest evidence is isotope labeling, introduced in Chapter 6: when people breathe oxygen labeled with oxygen-18, the label appears in newly produced water, not exhaled \chem{CO_2}. This also explains why the overall equation is misleading: the oxygen atoms on its two sides are not the same groups. Only tracing reveals their journeys.
\end{wuqu}

\section{Back to Those Thirty Squats}

We can now explain everything that happened during the opening minute. Your muscle cells consumed ATP furiously, and mitochondria ran flat out. The citric-acid cycle stripped electrons faster, the electron-transport chain carried them at full speed, and protons flooded ATP synthase. Demand at the electron exit surged, so breathing deepened and accelerated. Panting \textbf{replenished the terminal receiving stations} of countless electron-transport chains while removing the cycle's \chem{CO_2}.

What about heat? Electron-transport coupling is not perfect: each round loses some free energy directly as heat. Muscle contraction itself converts only about a quarter of its energy into mechanical work. With tens of trillions of mitochondria running together, waste heat triggers temperature-control alarms: blood vessels dilate, your face reddens, and you sweat. The mist on your mirror contains some water just made at the chain's end. Oxygen you inhaled is now leaving as water.

Time to check: oxygen became water, while carbon dioxide came from food's skeletons. Did you guess correctly?

\section{Poisons and Slimming Drugs: A Vulnerable Chain}

Understanding a machine means understanding how it fails. Cyanide, \chem{CN^-}, is highly poisonous because the cyanide ion fits into complex IV's heme-iron center and blocks oxygen. Electrons cannot reach the end, so the chain stops. Strangely, victims can have oxygen-rich blood while their cells suffocate in abundance; their skin may appear unusually flushed because the blood's oxygen cannot be used. Carbon monoxide attacks differently: it occupies hemoglobin, as in Chapter 48, preventing oxygen delivery to tissues. One poison blocks the destination, the other the transport route. Both strike the same vulnerability: electrons need an exit.

The reverse lesson is even more alarming. In the 1930s, people discovered that the proton smuggler 2,4-dinitrophenol turned respiration's energy into heat and sold it as a slimming drug. Fat burned rapidly, but a slightly excessive dose drove body temperature uncontrollably upward, and many users died of hyperthermia. It is one of pharmacology's darkest chapters: a correct mechanism but a lethal dose window. Remember this book's recurring red line---\textbf{dosage, strength of evidence, and individual differences}---whenever advertisements promise metabolic miracle drugs. Think of this electron-transport chain first.

\textbf{Translator safety note:} DNP is a dangerous toxic chemical, not a weight-loss treatment to try at any dose. Do not buy or take it. The historical account does not establish a safe dosing window. See the \href{https://www.food.gov.uk/sites/default/files/media/document/food-crime-strategic-assessment-2020.pdf}{Food Standards Agency assessment of DNP hazards}.

Medicine also uses the chain beneficially. One action of the common glucose-lowering drug metformin is mild inhibition of complex I, subtly changing cellular energy status and thereby regulating metabolism. Many pesticides and antibiotics target variants of electron-transport chains too. Evolution, meanwhile, does not regard heat as waste: warm-blooded animals rely heavily on mitochondrial heat to maintain temperature. Your ability to survive winter without a down jacket owes something to the chain's imperfect efficiency.

One final preview: here, electrons flow from food to oxygen and release energy. Could the process run the other way, using outside energy to pull electrons from water, push them uphill, and attach them to carbon dioxide to build sugar? That requires a machine converting light into electron flow and a way to handle oxygen released by splitting water. Plants have both. Their power stations' core---a proton gradient plus ATP synthase---is almost identical to mitochondria's, and even their evidence stories echo each other. In Chapter 54, we visit chloroplasts to examine this reverse-running machine and ask how your exhaled water and leaves' released oxygen connect within one chemical cycle.

\begin{tiaozhan}
You may skip this box without losing the main argument. Proton-motive force has two components: electrical potential difference $\Delta\psi$ and concentration difference $\Delta\mathrm{pH}$, combined as $\Delta p = \Delta\psi - 59\,\mathrm{mV}\cdot\Delta\mathrm{pH}$. Its mitochondrial total is about \ud{220}{mV}. Allocate this energy to electrons: one NADH provides 2 electrons, and the three pumps move about 10 protons across the membrane. ATP synthase lets roughly 3 protons return per ATP made, while importing ADP and phosphate into the matrix costs roughly another proton's gradient. Thus one NADH is worth about $10/4 = 2.5$ ATP, the origin of the exchange rate above. \chem{FADH_2} enters at complex II, bypassing the first pump and moving 4 fewer protons, so it yields about 1.5 ATP. Check efficiency again: burning 1 mole of glucose releases about \ud{2870}{kJ}; hydrolyzing 30 moles of ATP releases about $30\times 30.5 \approx 915$ \ud{}{kJ}, storing roughly a third. Where do the other two thirds go? Your warmth in winter is the answer.
\end{tiaozhan}

\section*{Practice}

\begin{sikaoti}
\item Draw paired material- and energy-flow diagrams following a glucose molecule's carbon to \chem{CO_2}, electrons to \chem{O_2}, and energy to ATP and heat. Label each route's stations: cytoplasm, mitochondrial matrix, inner membrane, and intermembrane space.
\item A cyanide-poisoned person can have oxygen-rich blood while cells suffocate and skin becomes unusually flushed. Explain the immediate cause of each observation using electron-transport terminology.
\item Suppose a chemical makes many proton-permeable holes in the inner mitochondrial membrane. Predict: (a) does electron transport speed up, slow down, or remain unchanged? (b) What happens to ATP production? (c) What happens to body temperature? Can this reasoning explain deaths from DNP slimming drugs?
\item A volunteer inhales oxygen-18-labeled \chem{^{18}O_2}. Will the label \textbf{first} appear in large amounts in exhaled \chem{CO_2}, or in body water and exhaled water vapor? Explain and identify the popular claim this experiment overturns.
\item Bacteria lack mitochondria but can perform aerobic respiration. Where are their electron-transport chains and ATP synthase located? How does this relate to the hypothesis that mitochondria originated from bacteria?
\item Estimate: an adult at rest synthesizes about \ud{50}{kg} of ATP daily. Its molar mass is about \ud{507}{g\cdot mol^{-1}}. How many moles and molecules do your cells synthesize daily, using \ud{6.02\times 10^{23}}{mol^{-1}}? Given this enormous output, why are oral ATP supplements for healthy people essentially a waste of money?
\end{sikaoti}
